% Generated by tools/build_v5.py; edit the assembly rules, not this file.
\section{Introduction and contribution}
Estimating the fraction of a population lying above a latent threshold is
not the same task as classifying every member of that population. Observations
can be concentrated on promising or ambiguous members, and biased individual
classifications can be combined through telescoping estimators. These
possibilities make algorithm-independent lower bounds important.

We study the following basic observation model. A fresh coin has a latent
probability $P$ drawn from a population distribution $G$. Conditional on $P$,
its observations are independent Bernoulli trials. An algorithm may adapt to
its entire history: it can create new coins, interleave observations of old
coins, and stop at a random time. Its target is
\begin{equation}\label{eq:target}
 \theta(G)=\Prob_G(P>1/2).
\end{equation}
One coin flip is one unit of work. The principal obstruction is the
population's mass near $1/2$, but a single power exponent need not describe
that mass accurately.

\paragraph{Distributional and distribution-blind optimality.}
The central object is the entire function
\begin{equation}\label{eq:modulus-intro}
 \omega_\mu(h)=\int\min\{1,h/d\}\,\mu(\mathrm d d),
 \qquad \mu=\mathcal L(|P-1/2|).
\end{equation}
The main theorem identifies this function as the minimax error law. Thus
irregular, discrete, and logarithmically modified margin distributions are
covered by one statement, rather than by separate exponent calculations.
The same algorithm attains the law for every magnitude distribution, rather
than requiring a distribution-dependent allocation. A second theorem
characterizes unbiased feasibility by an inverse moment and proves a sharp
universal lower constant for the variance--work product.
Writing $F_\mu(h)=\Prob(D\le h)$ makes the two parts of the modulus explicit:
\[
 \omega_\mu(h)=F_\mu(h)+h\E_\mu[D^{-1};D>h].
\]
The first term represents unresolved near-threshold mass; the second records
the accumulated contribution of the remaining difficulty scales.

The lower bound makes many difficulty scales jointly hard: each scale
contributes to the target separation, while every next observation has
uniformly bounded information. The upper bound is a fixed geometric
randomization of capped sequential classifiers. Its conditional second
moment and expected work are both controlled by the same truncated inverse
margin. This common conditional bound is what permits arbitrary margin laws.
An exact two-success inverse-binomial calibration then chooses the number
of production samples without estimating either their mean cost or variance.
The underlying inverse-sampling identity is classical~\cite{haldane}; its
role here is to eliminate all margin knowledge from the budget-indexed rule.

\subsection{Relationship to established work}
\paragraph{Nested simulation.}
Uniform inner-sampling analyses include Lee--Glynn~\cite{leeglynn} and
Gordy--Juneja~\cite{gordyjuneja}; Broadie, Du, and Moallemi~\cite{broadie}
also identify Lee's 1998 thesis as an early analysis and introduce sequential
inner allocation for loss-probability estimation.
Giles and Haji-Ali~\cite{gileshajiali} combine adaptive inner sampling and
multilevel Monte Carlo (MLMC) to attain
$O(\varepsilon^{-2}|\log\varepsilon|^2)$ work for thresholded nested expectations.
Haji-Ali, Spence, and Teckentrup~\cite{hajispence} develop a broader adaptive
MLMC framework for probabilities. We do not claim these methods or the
critical upper rate as new. Our observed-state construction in
\cref{prop:continuous} verifies Assumptions 2.1 and 2.4 of
\cite{gileshajiali} literally and proves the matching lower rate on an
unstructured subclass. It does not address models with additional smooth
cross-state structure or alternative information oracles.
Giles's survey~\cite{gilessurvey} discusses the high kurtosis of level
corrections for discontinuous functions; our confidence amplification uses
only second moments, not an empirical variance estimate or a normal approximation.

\paragraph{Adaptive coin mixtures and information bounds.}
Lee and Valiant~\cite{leevaliant} prove fully adaptive optimal bounds for
mixtures separated from $1/2$ by a fixed gap, including dependence on the
positive fraction. Our gap-class specialization agrees with their rate in
the balanced regime (\cref{cor:gap}). The lower bounds here are driven by
sign priors near $1/2$; this explains why a uniform local posterior inequality
and the stopped KL chain rule suffice. Their more refined rare-event
analysis is not reproduced by our supremum over all sign functions.
The change-of-measure step is related to the standard sequential information
inequality of Kaufmann, Capp\'e, and Garivier~\cite{kaufmann}; it is not a
new general principle. Lee--Valiant also optimize estimators for known
positive and negative conditional bias laws under a per-coin depth constraint.
Here the magnitude law is arbitrary, the sign function is unrestricted,
and our lower bounds have no per-coin depth constraint.

\paragraph{Randomization and unbiased estimation.}
Single-term multilevel debiasing is established methodology
\cite{rheeglynn,vihola}. Blanchet and Glynn~\cite{blanchetglynn} obtain
finite work-normalized variance under regularity such as local twice
differentiability of the outer function. Syed and Wang~\cite{syedwang}
study repeatedly nested expectations under suitable regularity conditions.
Our threshold indicator is discontinuous: the inverse-margin condition and
all-algorithm converse characterize what remains possible in this setting.
The inverse-binomial sampling identity used for cost calibration dates to
classical work of Haldane~\cite{haldane}; we give a self-contained proof and
do not claim that identity as new.

\paragraph{Other observation models and scope.}
Brennan, Vinayak, and Jamieson~\cite{brennan} estimate mass above effect-size
thresholds from noisy test statistics, rather than unrestricted sequential
revisiting. Chen, Naslidnyk, and Briol~\cite{chen} study nested expectations
with kernel quadrature and identify minimax lower bounds as a research
direction. Their smooth-kernel setting differs from ours.
The proposed contribution is the arbitrary-margin minimax characterization,
a distribution-blind optimal-rate algorithm, and the inverse-margin unbiased
converse in the precisely specified experiment. The original arbitrary-margin theorem is not an instance-optimal or
rare-event theorem. Section~\ref{sec:rare} treats a global rarity
constraint separately; no general solution to all nested-expectation
problems is claimed.


\paragraph{Rare events and the verification boundary.}
With a known upper bound $\theta\le\rho$, positive mass may concentrate
on the most difficult magnitudes. The power-margin rates below and
above exponent two are matched by different constructions. At the
fixed quadratic law, a scale-entropy cost contributes an extra
$\sqrt{1+\log(1+\log(1/\rho))}$ beyond the proposed variational modulus.
Section~\ref{sec:rare} states the current results together; the full
rare-event proofs appear in the appendices. This corrects the earlier
matching-$\Psi$ conjecture rather than assuming it.

Both directions of the original RMS law are proved in Lean for the actual
hidden-magnitude experiment, with coefficients $1/64$ and $512$ and one
upper algorithm independent of the gap law. The critical minimax lower
bound is also proved in that observable model. The untruncated 713 estimator
and variational optimizer have separate formal specifications. The
confidence characterization, unbiased converse, critical upper construction,
and explicit comparison with $\Psi$ retain paper-proof status.
Section~\ref{sec:verification} records the exact scope and links.
